\section{Limitations}\label{sec:limits}

\paragraph{Scope}
The study covers two English Wikipedia question-answering datasets, six open encoders up to 0.6B parameters, two 4-bit 7--8B readers, one prompt and $k=5$ for the registered transfer ratios; the exploratory ratios at other $k$ show that the size of the ratio depends on how many passages the reader sees. Results may differ for other domains, longer answers, larger or proprietary readers, other prompts or agentic pipelines. All timings come from one GPU type (Tesla T4) and one CPU configuration.

\paragraph{Corpus and labels}
Except for Phase~3, retrieval ran on 300,000-passage sub-corpora that contain every gold passage and favour the dense encoders relative to BM25. Phase~3 tested the NQ ratios on the full corpus for BM25, MiniLM-L6, bge-small and e5-base-v2 only; bge-base, bge-large, Qwen3-Embedding-0.6B and all of HotpotQA rest on the sub-corpus. NQ relevance labels are sparse. The answer-string check is lexical; the human audit agreed with it on most items but covers 100 misses of one encoder on the sub-corpus, labelled by one annotator, the author. The joined NQ set covers \ev{N.nq.joined} of \ev{N.nq.judged} judged queries and may differ in difficulty from the excluded queries. A Phase~1 check found Hit@5 differences between joined and unjoined queries from \ev{D1B.min} to \ev{D1B.max} across the seven retrievers, and every 95\% interval included zero.

\paragraph{Metrics and statistics}
EM is strict. We report F1 and answer containment; correctness was judged by hand only for the 100 audited misses. The transfer ratio is a ratio of means and becomes unstable when the retrieval gain is small; our reporting rule removes the worst cases, but the bge-small intervals on NQ still include zero. The ratio depends on the baseline encoder, and the share of the gain assigned to conditional EM depends on how the decomposition is ordered (Tables~\ref{tab:altbase} and~\ref{tab:share}). The decomposition, the query-level accounting, the Both@5 ranking and ratios, the BM25 baseline and the GEE interaction model are post hoc. The accounting groups are defined by retrieval outcomes, so they locate the EM difference but do not explain it. The registered pooled GEE model assumed one retriever effect across datasets, which the interaction test did not support. Answer quality under approximate search was not measured. Multiplicity is controlled within each family, not across all tests in the paper. Bootstrap intervals, F1, answer containment and the pooled model were checked by a rerun on a second computer, not by an independent implementation.

\paragraph{Context experiment}
The hard negatives are top-ranked BM25 passages without a gold label, and a third of the questions had at least one that contained an answer string. Excluding those questions made the measured drop larger, so this does not weaken P5, but the subset analysis is a sensitivity check on 333 questions. The experiment used one reader and 500 questions. The random-passage control found no detected difference from the gold passage alone; it does not establish equivalence.

\paragraph{Judge and annotation}
The LLM judge was validated against one annotator, the author. With $\kappa$ = \ev{JA.kappa} the judge was excluded under the registered rule, so this limitation does not affect the conclusions, but no faithfulness result is reported.

\paragraph{Three study phases}
The predictions were registered in dated, version-controlled design documents in our own repository, not in a public registry. The Qwen--NQ cell and part of the $k$ ablation were run before the predictions were registered, and the software image changed between Phases~1 and~2. P7 and P8 were registered after the sub-corpus results were known; they test whether those results replicate, not a new hypothesis. Within-cell comparisons share one environment. For comparisons across phases, \ev{CS.same} of \ev{CS.n} identical prompts gave identical outputs.
